\documentclass[11pt]{article}
\usepackage[letterpaper,margin=1in]{geometry}
\usepackage{amsmath,amssymb,amsthm}
\usepackage{graphicx}
\usepackage[hidelinks]{hyperref}
\usepackage{microtype}

\newtheorem{theorem}{Theorem}
\newtheorem{lemma}{Lemma}
\newcommand{\C}{\mathbb C}
\newcommand{\1}{\mathbf 1}

\title{A scalar operator-valued measure with no separable dilation\\
\large Full negative answer to Problem B of arXiv:1110.5833}
\author{Automated mathematics run \texttt{fa\_banach\_001}\\
Lane 5, agent \texttt{agent\_lane\_05}}
\date{9 August 2026}

\begin{document}
\maketitle

\begin{center}
\fbox{\parbox{0.91\textwidth}{\textbf{Status.}
Candidate full solution, likely valid, pending human review.  The construction
answers both parts of the source problem negatively.}}
\end{center}

\section{Source question}

Han, Larson, Liu, and Liu develop Banach-space dilations of
operator-valued measures in \cite{HLLL}.  Their Problem B appears in Section
2.4, printed page 37 (PDF page 41), after they prove separability for the
discrete framing case and before Proposition 2.38.

\begin{figure}[ht]
  \centering
  \includegraphics[width=\textwidth]{figures/open_problem_crop.png}
  \caption{Problem B in the source paper, rendered from
  \texttt{source\_paper.pdf}, PDF page 41.}
\end{figure}

The problem asks:
\begin{quote}
Assume that $(\Omega,\Sigma,E,B(X))$ is an operator-valued measure
system such that $X$ is a separable Banach space. Is the dilation space
$\alpha$ equipped with the minimal dilation norm always separable? If not,
does there exist a dilation space which is separable?
\end{quote}

We answer both questions negatively, already for the one-dimensional Banach
space $X=\C$ and a positive scalar multiple of the identity.

\section{Definitions}

For an operator-valued measure $E:\Sigma\to B(X)$, the source's algebraic
elementary dilation space is spanned by the set functions
\[
   E_{B,x}(A)=E(A\cap B)x
   \qquad (A,B\in\Sigma,\ x\in X).
\]
Its minimal dilation norm is
\begin{equation}\label{eq:minimal-norm}
 \left\|\sum_{k=1}^n c_k E_{B_k,x_k}\right\|_\alpha
   =\sup_{A\in\Sigma}
     \left\|\sum_{k=1}^n c_k E(A\cap B_k)x_k\right\|_X.
\end{equation}
The minimal dilation space is the completion in this norm.

A dilation of $E$ consists of a Banach space $Z$, bounded maps
$T:X\to Z$ and $S:Z\to X$, and a countably additive projection-valued
measure $F:\Sigma\to B(Z)$ such that
\begin{equation}\label{eq:dilation}
                 E(B)=S F(B)T \qquad(B\in\Sigma).
\end{equation}
Projection-valued measures in the source are spectral:
$F(A)F(B)=F(A\cap B)$.

\section{Proof intuition}

The obstruction is not the dimension of $X$ but the density of the underlying
measure algebra.  An uncountable Bernoulli product has uncountably many
coordinate events separated by a fixed amount of measure.  In the minimal
model, their elementary vectors inherit this separation directly from
\eqref{eq:minimal-norm}.

Changing the dilation cannot compress the obstruction into a separable space.
Indeed, every coordinate event produces a vector $F(B_i)T1$.  If two such
vectors were close, applying the spectral projection of the one-sided
difference $B_i\setminus B_j$ would isolate a vector whose scalar coefficient
is exactly $\mu(B_i\setminus B_j)=1/4$.  Uniform boundedness of the spectral
projections turns this into a fixed norm separation.

\section{Full negative solution}

\begin{theorem}\label{thm:main}
There exist a measurable space $(\Omega,\Sigma)$, a separable Banach space
$X$, and an operator-valued probability measure $E:\Sigma\to B(X)$ such that

\begin{enumerate}
\item the minimal dilation space of $E$ is nonseparable; and
\item every projection-valued dilation space of $E$ is nonseparable.
\end{enumerate}

Consequently both questions in Problem B of \cite{HLLL} have negative
answers.
\end{theorem}

\begin{proof}
Let $I$ be an uncountable set and take the fair Bernoulli product probability
space
\[
  (\Omega,\Sigma,\mu)
    =\bigl(\{0,1\}^{I},\ \text{product }\sigma\text{-algebra},\
                    (\tfrac12\delta_0+\tfrac12\delta_1)^I\bigr).
\]
Let $X=\C$ and define
\[
                    E(B)=\mu(B)I_\C \qquad(B\in\Sigma).
\]
This is a strongly countably additive $B(\C)$-valued probability measure, and
$X$ is separable.

For $i\in I$, put
\[
             B_i=\{\omega\in\Omega:\omega_i=1\},
             \qquad u_i=E_{B_i,1}.
\]
If $i\ne j$, independence of the two coordinates gives
\begin{equation}\label{eq:quarter}
             \mu(B_i\setminus B_j)=\frac14.
\end{equation}
Using $A=B_i\setminus B_j$ in \eqref{eq:minimal-norm}, we obtain
\[
\begin{split}
 \|u_i-u_j\|_\alpha
 &\geq
 \bigl|E(A\cap B_i)1-E(A\cap B_j)1\bigr|  \\
 &=\mu(B_i\setminus B_j)=\frac14.
\end{split}
\]
Thus the algebraic minimal dilation space already contains the uncountable
$1/4$-separated family $(u_i)_{i\in I}$.  Its completion is nonseparable.

It remains to show that no other dilation can be separable.  Let
$(Z,F,S,T)$ be an arbitrary dilation as in \eqref{eq:dilation}.  Since
$X=\C$, write $z=T1\in Z$ and regard $S$ as a functional in $Z^*$.  Then
\begin{equation}\label{eq:coefficient}
                    S(F(B)z)=\mu(B) \qquad(B\in\Sigma).
\end{equation}
The family of projections is uniformly bounded:
\begin{equation}\label{eq:M}
                       M:=\sup_{B\in\Sigma}\|F(B)\|<\infty.
\end{equation}
Indeed, for each $y\in Z$, the map $B\mapsto F(B)y$ is a countably additive
$Z$-valued measure, hence has bounded range; the uniform boundedness principle
then gives \eqref{eq:M}.

Set $w_i=F(B_i)z$.  Fix distinct $i,j$ and let
$D=B_i\setminus B_j$.  Spectrality gives
\[
 F(D)(w_i-w_j)
   =F(D\cap B_i)z-F(D\cap B_j)z
   =F(D)z.
\]
Together with \eqref{eq:coefficient}, \eqref{eq:quarter}, and
\eqref{eq:M}, this yields
\[
 \|w_i-w_j\|
 \geq \frac{\|F(D)z\|}{M}
 \geq \frac{|S(F(D)z)|}{M\|S\|}
 =\frac{1}{4M\|S\|}.
\]
Here $S\ne0$ because $S(F(\Omega)z)=\mu(\Omega)=1$.
Therefore $(w_i)_{i\in I}$ is an uncountable uniformly separated family in
$Z$, so $Z$ cannot be separable.  This proves both assertions.
\end{proof}

\section{Verification and scope}

The proof has no numerical or unproved dependency.  Its checkable points are:
\begin{enumerate}
\item $X=\C$ satisfies the separability hypothesis and $E(\Omega)=I_\C$;
\item distinct fair Bernoulli coordinates satisfy
      $\mu(B_i\setminus B_j)=1/4$;
\item the displayed minimal norm is exactly the source's definition;
\item strong countable additivity of $F$ implies pointwise boundedness of
      $\{F(B):B\in\Sigma\}$, and Banach--Steinhaus makes it uniformly bounded;
\item spectrality isolates $F(B_i\setminus B_j)z$ from $w_i-w_j$.
\end{enumerate}

The construction uses a non-countably-generated measure algebra.  It therefore
does not contradict Proposition 2.38 of the source, which obtains separability
from a countable Borel basis together with total strong countable additivity.
Rather, the example shows that some restriction of this kind is essential if
one wants a general separability theorem.

\section{Bounded novelty check}

On 9 August 2026 we searched the run's registry, solution, attempt, and
proof-gap indexes; the local arXiv corpus; exact phrases from Problem B; and
web searches combining ``minimal dilation space'', ``operator-valued
measure'', and ``separable dilation''.  The search found the source and the
authors' later exposition arXiv:1411.4559, but no claimed solution of Problem
B.  The later exposition repeats several other memoir problems while omitting
this one.  This is evidence, not a proof of novelty; the result should receive
expert literature review.

\section{Human-review recommendation}

Review as a candidate full negative answer.  The main verifier focus should be
the intended generality of Problem B (arbitrary measurable spaces) and the
standard bounded-range/uniform-boundedness step for $F$.  If both are accepted,
the argument closes the entire stated problem rather than merely a subcase.

\begin{thebibliography}{1}
\bibitem{HLLL}
D. Han, D. R. Larson, B. Liu, and R. Liu,
\emph{Operator-Valued Measures, Dilations, and the Theory of Frames},
arXiv:1110.5833 (2011); Memoirs of the American Mathematical Society
\textbf{229}, no. 1075.
\end{thebibliography}

\end{document}
